\section{操作性微观基线：构造、最小性与填充不变性}
\label{app:operational-baseline}

原始编码的状态数不是表示不变量：给每个状态附加一个不影响任何指定观测和
动力学的复制标签，就能任意放大它。本附录给出正文所用操作性微观空间的完整
构造。这个空间相对于预先声明的探针与允许动力学族是规范的；它不是脱离系统
边界的绝对``真实微观空间''。

\subsection{共同稳定的观测代数}

设 $X$ 为有限集，$u:X\twoheadrightarrow U$ 是在选择宏观表示之前固定的微观
探针，$\cK\subseteq\Stoch(X,X)$ 是非空允许动力学族。令
\[
  \cA^{\mathrm{mic}}_0:=u^*\mathbb R^U
\]
并递归定义
\begin{equation}
  \cA^{\mathrm{mic}}_{n+1}
  :=\operatorname{alg}\!\left(
      \cA^{\mathrm{mic}}_n
      \cup\bigcup_{K\in\cK}P_K\cA^{\mathrm{mic}}_n
    \right),
  \label{eq:app-micro-algebra-iteration}
\end{equation}
其中 $\operatorname{alg}$ 表示 $\mathbb R^X$ 内生成的最小含幺点乘子代数。

\begin{theorem}[操作性微观商]
\label{thm:app-operational-core}
序列~\eqref{eq:app-micro-algebra-iteration} 在有限步后稳定。设稳定代数为
$\cA_{\mathrm{mic}}$，并令其诱导分区的商映射为
\[
  p:X\twoheadrightarrow B.
\]
则：
\begin{enumerate}
  \item 存在唯一 $\bar u:B\to U$ 使 $u=\bar u\circ p$；
  \item 对每个 $K\in\cK$，存在唯一 $\bar K\in\Stoch(B,B)$ 使
  \begin{equation}
    KQ_p=Q_p\bar K;
    \label{eq:app-operational-intertwining}
  \end{equation}
  \item $p$ 是同时满足前两项的最粗确定性商：若
  $p':X\twoheadrightarrow B'$ 保留 $u$，并且每个 $K\in\cK$ 都在 $p'$
  上精确闭合，则存在唯一满射 $s:B'\twoheadrightarrow B$ 使
  $p=s\circ p'$。
\end{enumerate}
因此 $(B,\bar u,\{\bar K:K\in\cK\})$ 在典范同构意义下唯一。
\end{theorem}

\begin{proof}
代数序列单调增加。每次严格增加都会提高其作为 $\mathbb R^X$ 子空间的维数，
故至多有限次严格增加；一旦相邻两项相等，式~\eqref{eq:app-micro-algebra-iteration}
保证其后各项也相等。在稳定处，
\[
  P_K\cA_{\mathrm{mic}}\subseteq\cA_{\mathrm{mic}}
  \qquad(K\in\cK).
\]
有限观测代数与分区的对应由引理~\ref{lem:app-finite-algebra-partition} 给出：
$\cA_{\mathrm{mic}}=p^*\mathbb R^B$。于是
定理~\ref{thm:app-four-equivalences} 的观测闭合条件对每个 $K$ 成立，逐一给出
唯一 $\bar K$ 与式~\eqref{eq:app-operational-intertwining}。又因
$u^*\mathbb R^U\subseteq\cA_{\mathrm{mic}}$，$u$ 在 $p$ 的每个纤维上
常值，故唯一分解为 $u=\bar u\circ p$。

现设 $p'$ 满足题设条件。其分区代数
$\cA_{p'}:=(p')^*\mathbb R^{B'}$ 包含 $u^*\mathbb R^U$，并对全部
$P_K$ 不变。对 $n$ 归纳可得
$\cA^{\mathrm{mic}}_n\subseteq\cA_{p'}$，从而
$\cA_{\mathrm{mic}}\subseteq\cA_{p'}$。代数包含方向与分区细化方向相反，
所以 $p'$ 的每个纤维包含于某个 $p$ 的纤维；即 $p$ 在 $p'$ 上常值。
因此存在唯一 $s:B'\twoheadrightarrow B$ 使 $p=s\circ p'$。满射性来自
$p,p'$ 的满射性。
\end{proof}

\subsection{等价的稳定分区细化}

令 $\Pi_0$ 为 $u$ 的纤维分区。已知 $\Pi_n$ 后，定义 $\Pi_{n+1}$ 所对应的
等价关系为
\begin{equation}
 x\equiv_{n+1}x'
 \quad\Longleftrightarrow\quad
 \begin{cases}
   x\equiv_nx',\\
   K(x,C)=K(x',C),&
   \forall K\in\cK,\ \forall C\in\Pi_n.
 \end{cases}
 \label{eq:app-stable-refinement}
\end{equation}

\begin{proposition}[稳定分区算法]
\label{prop:app-stable-refinement}
$\Pi_n$ 正是 $\cA_n^{\mathrm{mic}}$ 的分区。序列在有限步后稳定，稳定分区
就是 $p:X\twoheadrightarrow B$ 的纤维分区。它是所有细化 $u$ 的纤维分区、
并对每个 $K\in\cK$ strong-lumpable 的分区中最粗者。
\end{proposition}

\begin{proof}
若 $\Pi_n$ 对应 $\cA_n^{\mathrm{mic}}$，则该代数由块指示函数
$\{\one_C:C\in\Pi_n\}$ 张成，且
\[
  (P_K\one_C)(x)=K(x,C).
\]
生成代数所诱导的等价关系恰要求旧块标签相同，并要求所有上述函数值相同；
这就是式~\eqref{eq:app-stable-refinement}。由 $n=0$ 归纳得到第一项，终止性
与定理~\ref{thm:app-operational-core} 相同。稳定时，块转移概率在每个块内
对每个 $K$ 恒定，故由定理~\ref{thm:app-four-equivalences} 得到共同
strong lumpability。反之，任一满足题设的稳定分区代数都包含初始代数且对
全部 $P_K$ 不变，定理~\ref{thm:app-operational-core} 的最小性证明说明它
必细化 $p$ 的分区。
\end{proof}

当 $\cK$ 是有限且各转移概率可比较时，式~\eqref{eq:app-stable-refinement}
给出直接有限算法；对抽象无限核族，它仍是数学刻画，但不自动给出有限输入
意义下的可执行检验程序。

若任务 $g:X\twoheadrightarrow O$ 由已声明微观信息确定，即 $g=t\circ u$，
则存在唯一 $\bar g:B\twoheadrightarrow O$ 使 $g=\bar g\circ p$。正文所称
宏观表示是 $B$ 的任务保真商
\[
  s:B\twoheadrightarrow Y,
  \qquad \bar g=r\circ s.
\]
只有在 $|Y|<|B|$ 时，$s$ 才是严格信息商。状态比 $|Y|/|B|$ 在这个信息
关系固定之后才具有本文所需的操作含义。

\begin{proposition}[约化前后闭合缺陷相同]
\label{prop:app-operational-defect-equality}
对任意 $K\in\cK$、其诱导核 $\bar K$、任务保真商
$s:B\twoheadrightarrow Y$ 以及任意 $L\in\Stoch(Y,Y)$，令
$q=s\circ p$，则
\begin{equation}
 d_\infty(KQ_q,Q_qL)
 =
 d_\infty(\bar KQ_s,Q_sL).
 \label{eq:app-operational-defect-equality}
\end{equation}
\end{proposition}

\begin{proof}
由 $Q_q=Q_pQ_s$ 和式~\eqref{eq:app-operational-intertwining}，
\[
 KQ_q=Q_p\bar KQ_s,
 \qquad
 Q_qL=Q_pQ_sL.
\]
$p$ 满射，故引理~\ref{lem:app-surjective-pullback} 表明前复合 $Q_p$
保持行最大距离，得到式~\eqref{eq:app-operational-defect-equality}。
\end{proof}

\subsection{精确复制填充不变性}

写 $[M]:=\{1,\ldots,M\}$。设原动力学族以同一索引集表示为
$\cK=\{K_\theta:\theta\in\Theta\}$。考虑
\[
 X^{(M)}:=X\times[M],
 \qquad
 u^{(M)}(x,i):=u(x),
 \qquad
 g^{(M)}(x,i):=g(x),
\]
以及核族 $\cK^{(M)}=\{K_\theta^{(M)}:\theta\in\Theta\}$，并假设
\begin{equation}
 \sum_{j=1}^M
 K_\theta^{(M)}\bigl((x,i),(x',j)\bigr)
 =K_\theta(x,x')
 \quad
 (\theta,x,x',i).
 \label{eq:app-exact-padding}
\end{equation}

\begin{theorem}[无意义复制填充不变性]
\label{thm:app-padding-invariance}
在条件~\eqref{eq:app-exact-padding} 下，填充系统的操作性微观空间
$B^{(M)}$ 与 $B$ 典范同构；该同构下，对应索引 $\theta$ 的诱导核相同。
因而任一任务保真商 $s:B\twoheadrightarrow Y$ 的状态比 $|Y|/|B|$、精确
闭合性质及正文定义的最坏行闭合缺陷均不随 $M$ 改变。
\end{theorem}

\begin{proof}
对稳定分区细化做归纳。初始分区的块具有
$C\times[M]$ 的形式，其中 $C$ 是 $u$ 的一个纤维。若 $\Pi_n^{(M)}$
由原分区 $\Pi_n$ 的块如此提升，则对任意 $C\in\Pi_n$，
\begin{align*}
 K_\theta^{(M)}\bigl((x,i),C\times[M]\bigr)
 &=\sum_{x'\in C}\sum_{j=1}^M
   K_\theta^{(M)}\bigl((x,i),(x',j)\bigr)\\
 &=\sum_{x'\in C}K_\theta(x,x')
 =K_\theta(x,C).
\end{align*}
该值与复制标签 $i$ 无关，而且两个原状态在下一轮是否被分开，恰由原系统
的对应块转移概率决定。因此
\[
  \Pi_n^{(M)}=\{C\times[M]:C\in\Pi_n\}
\]
对每个 $n$ 成立。稳定后，$C\times[M]\mapsto C$ 给出
$B^{(M)}\cong B$。

对两个稳定块 $C,D$，填充后诱导核从 $C\times[M]$ 到 $D\times[M]$
的概率由上式等于原诱导核从 $C$ 到 $D$ 的概率，所以对应诱导核相同。
任意后续商 $s$ 的状态数不变；其交换式两边在此同构下逐行相同，故精确性质
与 $d_\infty$ 缺陷也不变。
\end{proof}

定理~\ref{thm:app-padding-invariance} 只排除满足
式~\eqref{eq:app-exact-padding} 的\emph{精确}无意义标签。若标签对探针或有限
时间动力学有小但非零影响，是否把它视为可忽略自由度取决于另行声明的测量
分辨率、统计置信度与时间尺度；本定理没有给出与这些选择无关的阈值。
